\paragraph{Cache replacement policies}
This is mostly a recap from SPCA, but again, the DB can be more aggressive at optimizing, because it has much more information on the data accesses.

A commonly used strategy is Least Recently Used (LRU).
A buffer is evicted, when space is needed and it is the least recently used buffer in the buffer pool.

As much sense as this policy makes in operating systems, in DBMS, it doesn't really work (but is still used in some systems).
There are a number of possible issues with LRU:
\begin{itemize}
    \item \bi{Table Scan Flooding}: If a large table is loaded to be scanned just once, it will flush out the entire cache, also pages that are used \textit{very often}.
    \item \bi{Index Range Scan}: A range scan using an index will pollute the cache with random pages
\end{itemize}
